\clearpage
\section{Repeated Scalar Execution}\label{page:repeated-scalar-execution}

This section defines the architectural role of repeated scalar execution in Bedrock. The architecture exposes
scalar register state and repeated-instruction capabilities.

\subsection{Architectural Scalar Semantics}
REP, REPcc, REPG, and any explicitly defined repeat instruction are architecturally scalar. Their visible behavior is
the same as executing the repeated instruction or byte-counted instruction group in program order under the selected
counter and termination rules.
The architectural counter, FLAGS image, memory effects, and auto-update register effects are observed as if each
committed iteration completed before the next one began.

The body is decoded when repeat state is entered and remains fixed for that repeat context. Each iteration evaluates
register and memory operands by the body's ordinary rules, including the normal scalar effects of auto-update
effective addresses. FLAGS and FFLAGS follow the body's instruction semantics and any repeat-specific accumulation
rule. Faults remain precise at the repeated-operation boundary.

\subsection{Repeat Syntax and Body Eligibility}

The accepted spellings are:
\begin{itemize}
\item \texttt{REPcc Rn, (<instruction>)}
\item \texttt{REP Rn, (<instruction>)}
\item \texttt{REPG Rn, \{ <instruction>... \}}
\item \texttt{REPGF Rn, \{ <instruction>... \}}, an assembler-only spelling that emits REPG after applying the
additional eligibility checks described below
\end{itemize}

The processor decodes the body and checks its repeat eligibility before applying any zero-based termination rule. A
malformed or unavailable body therefore faults even when the selected register is zero. Once a valid body has
passed those checks, a zero captured count annuls REP or REPcc without architectural side effects. REPG has no entry
zero-count annul; it begins the group and applies its live loop-control rule only at a successful group boundary.

REPGF is an assembler-checked spelling that emits REPG after applying the additional body-eligibility checks. The
assembler accepts only candidate bodies that do not write the selected loop-control register. The emitted encoding,
decoded operation, and repeat-continuation state are identical to REPG.

\subsection{Counter, Condition, and Control Rules}

Entering REP or REPcc state captures the selected Rn value as an unsigned remaining count. Each successfully
committed iteration decrements that captured count. A condition-false REPcc iteration still commits its body and
performs the decrement before repeat state is cleared. A valid zero count performs no body side effects after
decoding and eligibility checks complete.

During REP and REPcc, a body read of the selected counter register observes the iteration-start architectural value,
while a body write to that register is ignored. REPG instead uses the selected architectural Rn as its live unsigned
loop-control value and captures no separate remaining count. REPG executes reads and writes of that register in
normal in-group program order. At each successful whole-group boundary, it examines the current value after all body
writes. If the value is zero, REPG leaves it unchanged and exits. Otherwise, REPG first decrements it by one and then
uses the resulting value to select the next action: zero exits, while nonzero continues at the group start. This
guarded decrement and resulting-value decision form the group-boundary state transition. When the body does not
write the register, an initial zero executes one iteration and an initial positive value N executes exactly N
iterations; both cases leave zero. Any body write controls the boundary reached by that iteration, so writing either
zero or one terminates at that boundary. The checked REPGF spelling rejects such a write at assembly time.

The \texttt{REPcc observation} attribute in every REPcc-capable instruction description specifies the value used for
the condition test. A \texttt{result:operand} or \texttt{source:operand} observation names an actual operand. A
\texttt{computed} observation names a derived value defined by that instruction: CMP observes its selected-width
temporary difference, TEST observes its selected-width temporary conjunction, BTEST/BSET/BCLR/BCHG observe the old
tested bit as a B-sized 0 or 1, the signed and unsigned bounds checks observe their violation result as a B-sized 0 or
1, and SEGLEA observes its segment-bound failure result as a B-sized 0 or 1.

REPcc constructs a temporary condition image from that value: Z is one exactly when the observed value is zero, N is
its most significant bit at the observation width, and C and V are zero. The condition test never samples
architectural FLAGS, and the REPcc wrapper does not write architectural FLAGS; any FLAGS effects of the body itself
still commit normally. The observed value is captured before repeat counter-destination suppression. The first
condition test occurs after the first body iteration. Body completion, value observation, temporary condition-image
construction, body architectural commit, counter decrement, and continuation-PC selection form one precise atomic
commit step. Consequently, an interrupt handler cannot change the decision for an iteration that has reached this
step. REP is the unconditional T-condition spelling, and F is reserved for REPcc.

A repeat body may not contain another repeat operation, a repeat-control instruction, or a control transfer.
While repetition is active, PC names the current body instruction, and the initially decoded instruction or group
remains fixed. Handler entry saves repeat state in FRAME\_EXT1 and clears the active architectural repeat state; ERET
restores that state atomically with the saved PC and control state. TF and RF apply to the trace units defined by the
Event Priority and Debug Trace subsection.

\subsection{Fault and Restart Semantics}
A fault during repeated execution preserves committed work and saves the state needed to resume or emulate the
remainder. Event entry records the active repeat context in FRAME\_EXT1 and clears active repeat state. The Repeat
Continuation subsection of the Architectural Event Processing Model defines the saved format and ERET restoration.

REP and REPcc atomically commit the one-instruction body, repeat-observed value and temporary condition, counter decrement, and
continuation-PC selection. If the body faults, that iteration does not commit and the fault is reported at the body
instruction. Every committed iteration decrements the counter, including a condition-false iteration that terminates
REPcc; earlier iterations remain committed.

For REPG, the unsigned body byte count must describe a group that ends on an instruction boundary. Instructions in
the group commit in program order. If an instruction faults, earlier completed instructions in the current group
iteration remain committed and later instructions do not execute. The fault is reported at the faulting instruction;
the live loop-control register retains every committed body effect, but the group-boundary state transition does not
occur for the incomplete group iteration. The saved PC and repeat continuation retain the information required to
reconstruct the group and continue the remaining work.
